% q0013.tex — Two Intersections (Simple vs Compound Interest)
\question{
  id        = {0013},
  title     = {Two Intersections},
  source    = {OBECON},
  year      = {2026},
  round     = {Seletiva IEO -- Phase A},
  language  = {English},
  topic     = {Finance},
  subtopic  = {Time Value of Money / Interest},
  type      = {objective},
  statement = {%
    Consider two investment options:
    \begin{enumerate}
      \item An investment that offers \textbf{simple interest} at rate $r$:
            $A_s = P(1 + rt)$
      \item An investment that offers \textbf{compound interest} at rate $r$:
            $A_c = P(1+r)^t$
    \end{enumerate}
    where $A$ is the accumulated amount, $P$ is the principal, $r$ is the
    annual rate (decimal), and $t$ is the time in years.

    Carlos invests \$1,000 in the compound option; Frederico invests \$1,000
    in the simple option. What can we assume about their decisions?
  },
  choices   = {%
    \item Carlos made the optimal choice, as compound interest always yields
          higher returns than simple interest.
    \item Frederico made the optimal choice, as simple interest is more
          predictable and less risky than compound interest.
    \item Both Carlos and Frederico may have made the optimal choice,
          depending on the investment duration.
    \item Both will receive the same amount over the long term, as the
          difference between simple and compound interest diminishes over time.
    \item Both options will yield the same return after $\dfrac{P}{r}$ years.
  },
  answer    = {C},
  solution  = {%
    Compare $A_s = P(1+rt)$ and $A_c = P(1+r)^t$:
    \begin{itemize}
      \item At $t = 0$: both equal $P$.
      \item At $t = 1$: $A_s = P(1+r) = A_c$. \textbf{Equal.}
      \item For $0 < t < 1$: $(1+r)^t < 1 + rt$ by concavity of power
            function~$\Rightarrow$ \textbf{simple interest is higher}.
            (Frederico does better for short-term investments.)
      \item For $t > 1$: $(1+r)^t > 1 + rt$~$\Rightarrow$
            \textbf{compound interest is higher}.
            (Carlos does better for long-term investments.)
    \end{itemize}

    Therefore, neither choice is universally optimal; it depends on the
    investment duration. Both may have made the correct decision.
    Answer: \textbf{(C)}.
  },
}
